\section{Conclusion}
On a à présent calculé le groupe fondamental du Tore, du plan projectif réel et de la bouteille de Klein (à isomorphisme près). Du côté du Tore, on le visualise correctement car il se plonge dans $\mathbb{R}^3$ sans difficulté. Ce n'est pas le cas du plan projectif réel, ni de la bouteille de Klein, dont le plongement en $3$ dimensions induit une auto-intersection.\medskip\\
Visuellement, il n'est pas toujours évident de s'imaginer la transformation à partir d'un simple carré en nos surfaces (sauf pour le Tore, là c'est plutôt direct). Dans le but de montrer cette transformation, un programme codé en Python est à disposition dans le projet.\medskip\\
En le lançant, on a le choix d'observer un gif simulant la transformation du carré en la surface de notre choix (avec pour la bouteille de Klein, la possibilité de l'observer sous $2$ formes différentes, le $8$ et le vase), ou de générer un espace en 3d dans lequel on peut tourner la surface et l'observer.\medskip\\
Pour conclure ce document, voici une modeste bibliographie des livres que j'ai pu utiliser pour m'aider dans ce projet :
\addcontentsline{toc}{section}{Bibliographie}
\begin{thebibliography}{9}

\bibitem{hatcher}
Allen Hatcher,
\textit{Algebraic Topology},
Cambridge University Press, 2001.

\bibitem{engelking}
Ryszard Engelking,
\textit{General Topology},
Heldermann Verlag, 2e édition, 1989.

\bibitem{zisman}
Michel Zisman,
\textit{Topologie Algébrique Élémentaire},
Armand Colin, 1972.

\bibitem{munkres}
James R. Munkres,
\textit{Topology},
Prentice Hall, 2e édition, 2000.

\end{thebibliography}